% Auto-generated by paper/tables/gen_tab7_vllm.py
\begin{table}[t]
\centering
\renewcommand{\arraystretch}{1.12}
\setlength{\tabcolsep}{2pt}
\scriptsize
\begin{threeparttable}
\caption{Measured \spear gain under vLLM production-stack deployment.\tnote{a}}
\label{tab:vllm}
\begin{tabular*}{\columnwidth}{@{\extracolsep{\fill}}llllrrr@{}}
\toprule
\textbf{Case} & \textbf{Model} & \textbf{Pool} &
\textbf{Gap}\tnote{b} &
\textbf{$\Delta$HR} & \textbf{$\Delta$TTFT}\tnote{c} &
\textbf{$\Delta$e2e} \\
 & & & & (pp) & (\%) & (\%) \\
\midrule
Cache-bound & 1.5B & 60K & Tight & \textbf{+3.4} & $-35$ & $-72$ \\
Capacity-stressed & 1.5B & 30K & Loose & \textbf{+2.5} & $-6$ & $-4$ \\
Native-model & 30B & 60K & Tight & \textbf{+1.6} & $-1$ & $-3$ \\
\bottomrule
\end{tabular*}
\begin{tablenotes}\scriptsize
\item[a] Workload is GAIA/Qwen3 at 32-way concurrency.
1.5B is Qwen2.5-1.5B and 30B is the trajectory-native
Qwen3-Coder-30B-AWQ. Pool sizes are 60K and 30K tokens.
\item[b] Tight/Loose denotes the diverse tool-gap schedule.
\item[c] Deltas are versus paired LRU in vLLM and
TTFT is the mean and e2e is p99, in percent.
\end{tablenotes}
\end{threeparttable}
\end{table}
